\documentclass[11pt]{article}
\usepackage[margin=1in]{geometry}
\usepackage[T1]{fontenc}
\usepackage{lmodern,amsmath,amssymb,amsthm,booktabs,longtable,array}
\usepackage[colorlinks=true,linkcolor=blue,urlcolor=blue,citecolor=blue]{hyperref}
\usepackage{microtype}
\newtheorem{theorem}{Theorem}
\newtheorem{lemma}[theorem]{Lemma}
\newtheorem{corollary}[theorem]{Corollary}
\newcommand{\Z}{\mathbb Z}
\newcommand{\N}{\mathbb N}
\newcommand{\Ham}{\operatorname{Ham}}
\newcommand{\verificationstatus}{The proof sources were verified on 26 September 2026 using Lean
v4.33.0-rc1 and pinned dependencies. This public revision preserves those
sources; the original successful build log is included in the repository.}
\title{Bilipschitz Equivalence of Integer Power Word Metrics\\
\large Proof and Lean Formalization Report}
\author{Qiao Li}
\date{4 November 2026}
\begin{document}
\maketitle

\begin{abstract}
For an integer $a\ge2$, let $d_a$ be the word metric on $\Z$ generated by
$\{\pm a^k:k\in\N\}$. We give a proof that all these metric spaces admit
bilipschitz bijections. The central construction is a relative Hamming
recoding that preserves two distinguished symbols, including their positions,
and commutes with changes between those symbols. This property turns an
arbitrarily long carry into the difference of two generators. The Lean
implementation uses a Hamming constant of $35$ and yields forward and reverse
constants $564$ and $321$ for bases $2$ and $3$.
\par\smallskip\noindent\textbf{Verification status:} \verificationstatus
\end{abstract}

\section{The question and precise result}
Define
\[
 \ell_a(z)=\min\left\{r\in\N:
 z=\sum_{j=1}^{r}\varepsilon_j a^{k_j},\quad
 \varepsilon_j\in\{-1,1\},\ k_j\in\N\right\},
 \qquad d_a(x,y)=\ell_a(y-x).
\]
The empty sum represents zero. The minimum exists because $1=a^0$ is a
generator. The question is whether $(\Z,d_p)$ and $(\Z,d_q)$ are
quasi-isometric for primes $p,q$, especially $p=2,q=3$; the map may be an
arbitrary map of metric spaces. The definitions and related questions appear
in Nathanson~\cite[Section 4, Problems 1--2]{nathanson}.

\begin{theorem}\label{thm:main}
For every $a,b\in\N$ with $a,b\ge2$, there exists a bijection
$f:\Z\to\Z$ such that, for all $x,y\in\Z$,
\begin{align}
 d_b(f(x),f(y))&\le a\bigl(2+35(b^2-1)\bigr)d_a(x,y),\label{eq:forward}\\
 d_a(x,y)&\le b\bigl(2+35(a^2-1)\bigr)d_b(f(x),f(y)).\label{eq:reverse}
\end{align}
Consequently these spaces are bilipschitz equivalent, and hence
quasi-isometric.
\end{theorem}

\begin{corollary}\label{cor:23}
There exists a bijection $f:\Z\to\Z$ satisfying
\[
 \frac{1}{321}d_2(x,y)\le d_3(f(x),f(y))\le564d_2(x,y)
 \quad(x,y\in\Z).
\]
In particular, a common bilipschitz constant is $564$.
\end{corollary}

The formal theorem is an existence theorem in classical mathematics.
It makes no assertion that the chosen bijection is computable, additive,
affine, finite-state, or continuous for an adic topology. The constants above are those proved in the formal development.

\section{Relative recoding of finite words}
Let $A$ be a countable alphabet with distinguished element $0$ and at least
two elements. Add two new symbols $P,M$, and let $\mathcal S(A)$ be the set
of finite-support functions $\N\to A\sqcup\{P,M\}$, based at $0$.
Equip it with Hamming distance. Write
\[
 F(x)=\{j:x_j\in\{P,M\}\}
\]
for the finite frozen set. A flip on a finite set of coordinates interchanges
$P$ and $M$ there and fixes all ordinary letters.

\begin{lemma}[Relative Hamming equivalence]\label{lem:relative}
For any two such alphabets $A,B$, there is a bijection
$T:\mathcal S(A)\to\mathcal S(B)$ with both $T$ and $T^{-1}$
$35$-Lipschitz for Hamming distance. At every coordinate it preserves
whether the letter is $P$, $M$, or ordinary. It also commutes with every
finite coordinate flip.
\end{lemma}

Ordinary Hamming equivalence over finite alphabets is established in
Domat--Hoganson--Lyman~\cite[Section 4.2]{dhl}. Their absorption and infinite
partition method motivates the construction below. The additional frozen
set and flip conditions are established explicitly here; the ordinary
Hamming theorem alone does not supply the carry argument.

\subsection{An absorption matching on an infinite set}
Let $E$ be a countably infinite set and let $A^{(E)}$ denote its finitely
supported $A$-valued functions. For a pointed countable alphabet $B$, relate
$x\in A^{(E)}$ to $(b,z)\in B\times A^{(E)}$ when $x,z$ differ in at
most one coordinate. Every vertex on each side has infinitely many
neighbors: choose a fixed nonzero letter of $A$ and insert it outside the
finite support. Enumerating either countable side and choosing a fresh
neighbor gives injections in both directions which respect this same
relation. The relation-preserving Schr\"oder--Bernstein theorem then gives
a bijection
\[
 \psi_E:A^{(E)}\longrightarrow B\times A^{(E)},
 \qquad \psi_E(x)=(b,z),\qquad \Ham(x,z)\le1.
\]
Reserve the edge $0\leftrightarrow(0,0)$ before the matching construction
to ensure $\psi_E(0)=(0,0)$. The inverse has the same displacement
property. Choose $\psi_E$ as a function of the set $E$ alone. This last
condition is essential when comparing different frozen sets.

\subsection{Assembly with a variable frozen set}
Fix a partition $\N=\bigsqcup_{i\in\N}B_i$ into infinite sets; formally,
$B_i$ is the fiber of the first component of the standard pairing inverse.
For a frozen set $F$, put $E_i=B_i\setminus F$, which remains infinite.
If $i\notin F$, apply $\psi_{E_i}$ to the ordinary input data on $E_i$,
retaining its modified $A$-word and storing its extra $B$-letter at
coordinate $i$. If $i\in F$, use the identity on $E_i$ and store no extra
letter. Copy every frozen label unchanged. At each editable coordinate,
pair the retained $A$-letter with its stored $B$-letter. This defines
\[
 \Phi_{A,B}:\mathcal S(A)\longrightarrow\mathcal S(A\times B).
\]
It is bijective: from an output, recover $F$, split its ordinary pairs,
and apply the corresponding inverse matching in each active block.
Only finitely many input blocks are nonzero, and every local matching fixes
zero; thus both directions preserve finite support. The construction
depends on the frozen positions but not on the values $P$ versus $M$.
It therefore commutes with all flips.

\subsection{Uniform control across different frozen sets}
Here is the estimate that makes the construction global. Given two input
words, let $D$ be their difference set, $d=|D|$, and let $\pi(j)=i$ for
$j\in B_i$. Only indices in
\[
 I=\pi(D)\cup D,\qquad |I|\le2d,
\]
can have changed block input, editable set, or active status. For $i\notin I$,
both words have the same editable set and the same data on $B_i$, and
coordinate $i$ agrees. They therefore use exactly the same local matching
and produce identical retained data on $B_i$ and the same extra letter
at coordinate $i$.

Within an affected block the retained $A$-word differs from the original
at at most one position for each of the two inputs. The difference count
of the retained first components is consequently at most
$d+2|I|\le5d$. The extra components differ at at most $|I|\le2d$
coordinates. Hence $\Phi_{A,B}$ is $7$-Lipschitz. For its inverse the
same block argument gives $d+2|I|\le5d$: the extra letter has already
been included in the source word. Thus $\Phi_{A,B}^{-1}$ is
$5$-Lipschitz. This reasoning includes changes in membership of $F$,
rather than only pairs of words with one fixed frozen set.

Apply the construction with the alphabets reversed and use the coordinate
swap $A\times B\cong B\times A$. Then
\[
 T=\Phi_{B,A}^{-1}\circ\operatorname{swap}\circ\Phi_{A,B}
\]
has constant $5\cdot7=35$ in both directions and retains all frozen and
flip properties. This proves Lemma~\ref{lem:relative}.

\section{Signed expansions and the carry estimate}
For $A\ge4$, set $D_A=\{-1,0,1,\ldots,A-2\}$.
Every integer has a unique finite expansion
\begin{equation}\label{eq:digits}
 n=\sum_{i\ge0}d_i A^i,\qquad d_i\in D_A.
\end{equation}
Indeed, the next quotient is $q=\lfloor(n+1)/A\rfloor$ and the remainder
$d=n-Aq$ belongs to $D_A$. For $n\ne0$, $|q|<|n|$, so the process
terminates. For uniqueness, the least differing digit in two finite
expansions would give a nonzero difference divisible by $A$ with absolute
value at most $A-1$, which is impossible.

Identify $A-2$ with $P$ and $-1$ with $M$; the other $A-2$ digits form
the ordinary alphabet. Its size is at least two. Let
$E_A:\Z\to\mathcal S(\{0,\ldots,A-3\})$ be the resulting expansion
bijection. Given the relative recoding $T$ for bases $A,B\ge4$, define
\[
 f_{A,B}=E_B^{-1}\circ T\circ E_A.
\]

Fix a generator edge $n\mapsto n+A^k$. Starting at position $k$, let $t$
be the first digit below $A-2$; it exists because the word is eventually
zero. The addition changes the digits $k,\ldots,t-1$ from $P$ to $M$
and increments only the digit at $t$. This follows from
$(A-2)+1=-1+A$ and applies to negative integers as well.

First perform just those frozen flips, producing a word $u$. Its recoded
image differs from $T(E_A(n))$ by precisely the same flips. Decoding in
base $B$, the numerical change is
\begin{equation}\label{eq:telescoping}
 -(B-1)\sum_{i=k}^{t-1}B^i=B^k-B^t.
\end{equation}
Its $B$-word length is at most two, including the case $t=k$, when the
sum is empty. The word $u$ differs from $E_A(n+A^k)$ at only one position,
so their recoded words differ at at most $35$ positions. A change of one
target digit has coefficient of absolute value at most $B-1$, costing at
most $B-1$ signed $B$-generators. Therefore
\[
 d_B\bigl(f_{A,B}(n+A^k),f_{A,B}(n)\bigr)\le2+35(B-1).
\]
Reverse generator edges follow by symmetry. Summing over an actual
shortest signed-power word yields
\begin{align*}
 d_B(f_{A,B}(x),f_{A,B}(y))&\le[2+35(B-1)]d_A(x,y),\\
 d_A(x,y)&\le[2+35(A-1)]d_B(f_{A,B}(x),f_{A,B}(y)).
\end{align*}
The second inequality is the same argument for $T^{-1}$, which also
commutes with flips. The expansion maps themselves need no Hamming
Lipschitz bound; equation~\eqref{eq:telescoping} controls long carries.

\section{Returning to all bases}
For $a\ge2$ and $z\in\Z$,
\[
 \ell_a(z)\le\ell_{a^2}(z)\le a\ell_a(z).
\]
The first inequality is inclusion of generating sets. For the second,
an even power of $a$ is an $a^2$-generator and an odd power is a sum of
$a$ equal $a^2$-generators. The statement also holds for negative
generators. Take $A=a^2,B=b^2$ in the preceding construction and combine
these comparisons. This proves~\eqref{eq:forward}--\eqref{eq:reverse},
including the small bases $2$ and $3$.

\section{Formal statement and verification boundary}
The Lean development defines the predicate \texttt{WordBound} inductively
using exactly $r$ signed powers. The function \texttt{wordLength} takes the least attainable
$r$, and \texttt{wordDist a x y} is the resulting length of $y-x$.
Thus the encoded distance is the original word distance, rather than a
digit-weight surrogate. \texttt{PowerInt a} is a type synonym for $\Z$
with this distance cast to $\mathbb R$, equipped with a proved
\texttt{MetricSpace} instance.

The main public statements are in \texttt{lean/Nathanson/Main.lean}:
\begin{center}\small
\begin{tabular}{>{\raggedright\arraybackslash}p{0.48\linewidth}>{\raggedright\arraybackslash}p{0.43\linewidth}}
\toprule
Lean declaration (namespace \texttt{Nathanson}) & Mathematical content\\
\midrule
\texttt{relativeEquiv35} & Lemma~\ref{lem:relative}, both directions and flips\\
\texttt{all\_bases\_bilipschitz} & Theorem~\ref{thm:main}, integer distances\\
\texttt{all\_bases\_metric\_bilipschitz} & Standard \texttt{LipschitzWith} and \texttt{AntilipschitzWith} for a metric-space equivalence\\
\texttt{prime\_bases\_metric\_bilipschitz} & Specialization to any two primes\\
\texttt{two\_three\_bilipschitz} & Corollary~\ref{cor:23}, constants $564,321$\\
\texttt{two\_three\_metric\_bilipschitz} & The same constants in the metric-space interfaces\\
\bottomrule
\end{tabular}
\end{center}

The implementation is pinned to Lean \texttt{v4.33.0-rc1} and the following
mathlib commit:
\begin{center}\small\texttt{79d0395a1825a6264ad5d269e35e60537518955e}.\end{center}
The proof is organized into matching, absorption, block assembly,
cross-mask control, signed digits, carry transfer, and metric comparison.
It includes the infinite matching and the variable-mask argument themselves;
these are not imported as unproved mathematical assumptions. The local
matching uses classical choice. The checked transitive logical dependencies
are exactly Lean's \texttt{propext}, \texttt{Classical.choice}, and
\texttt{Quot.sound}; the corresponding output is retained in
\texttt{lean/build.log}.

\subsection*{Public repository and reproduction}
The public project is available at
\begin{center}\small
\url{https://github.com/KeeneLee/integer-power-word-metrics}.
\end{center}
The repository contains \texttt{README.md}, this report as
\texttt{report.pdf}, its source as \texttt{report.tex}, the Lean project in
\texttt{lean/}, and an MIT \texttt{LICENSE}. The 14 proof modules are in
\texttt{lean/Nathanson/}; \texttt{lean-toolchain}, \texttt{lakefile.toml},
and \texttt{lake-manifest.json} inside \texttt{lean/} pin the build environment.
With elan installed, run from the repository root:
\begin{verbatim}
cd lean
lake exe cache get
lake build
\end{verbatim}
The target declarations and their \texttt{\#print axioms} commands are in
\texttt{lean/Nathanson/Main.lean}. The retained September build log reports
\texttt{Build completed successfully (3013 jobs)}; the original build rebuilt
all 14 project modules and used the pinned upstream mathlib cache.
This November revision updates the public presentation and project paths,
without changing the Lean proof sources or claiming a new build.

\subsection*{Acknowledgment}
AI tools assisted with the proof and Lean formalization.

\begin{thebibliography}{9}
\bibitem{nathanson}
M. B. Nathanson, \emph{Bi-Lipschitz equivalent metrics on groups, and a problem
in additive number theory}, Portugaliae Mathematica \textbf{68} (2011),
191--203. \href{https://arxiv.org/abs/0902.3254}{arXiv:0902.3254}.
\bibitem{dhl}
G. Domat, H. Hoganson, and R. A. Lyman,
\emph{Coarse geometry of homeomorphism groups: Classifying countable Stone spaces},
version 2, 24 July 2026, Section 4.2, especially Lemma 4.7 and Theorem 4.8.
\href{https://arxiv.org/abs/2607.01196v2}{arXiv:2607.01196v2}.
\end{thebibliography}
\end{document}
